\section*{Capítulo \ref{ch:b2:precambrian-earth} --- A Terra pré-cambriana}

\begin{solution}{exo:b2:precambrian-earth:1}
(a) \omterm{def:b2:precambrian-earth:eons}{Hadeano} $4567 - 4031 = \qty{536}{Myr}$; \omterm{def:b2:precambrian-earth:eons}{Arqueano} $4031 - 2500 = \qty{1531}{Myr}$;
\omterm{def:b2:precambrian-earth:eons}{Proterozoico} $2500 - 538.8 = \qty{1961}{Myr}$. (b) $(4567 - 538.8)/4567 = 0.882$: cerca de sete
oitavos. (c) $4567 - 4404 = \qty{163}{Myr}$ para o zircão; $4567 - 4030 = \qty{537}{Myr}$ para
as rochas mais antigas.
\end{solution}

\begin{solution}{exo:b2:precambrian-earth:2}
(a) Os fatores $e^{\lambda t}$ com $\lambda = \ln2/T$ e $t = \qty{3}{Gyr}$ são 1.593
(\ce{^{238}U}), 19.18 (\ce{^{235}U}), 1.160 (\ce{^{232}Th}) e 5.292 (\ce{^{40}K}): as
contribuições eram 3.053, 1.573, 2.424 e \qty{4.329}{pW}, \qty{11.38}{pW} no total, contra
\qty{4.906}{pW} hoje. (b) $11.38/4.906 = 2.32$. (c) O tório-232 lidera hoje; o potássio-40 liderava
então.
\end{solution}

\begin{solution}{exo:b2:precambrian-earth:3}
(a) $1616 + 18.0\times(28 - 32.2) = \qty{1540}{\celsius}$. (b) $32.2 + (1500 - 1616)/18.0 =
\qty{25.8}{\percent}$. (c) Uma análise de rocha total não é a composição de um líquido: um
derrame de \omterm{def:b2:precambrian-earth:greenstone}{komatiito} pode conter cristais de olivina acumulados a partir do líquido, que acrescentam MgO. Seus
\qty{38}{\percent} estão além dos \qty{32}{\percent} dos líquidos komatiíticos mais ricos, e
além dos experimentos em que a reta se apoia.
\end{solution}

\begin{solution}{exo:b2:precambrian-earth:4}
Rochas mais antigas, \qty{4.03}{Ga}; primeiros estromatólitos, \qty{3.43}{Ga}; fim dos domos do Pilbara,
\qty{3.1}{Ga}; \omterm{def:b2:precambrian-earth:goe}{Grande Evento de Oxidação}, cerca de 2.4 a \qty{2.1}{Ga}; últimos \omterm{def:b2:precambrian-earth:greenstone}{komatiitos} comuns,
\qty{1.9}{Ga}; segunda subida do oxigênio, cerca de \qty{0.8}{Ga}; fim da Marinoana,
\qty{0.635}{Ga}; primeiros fósseis ediacaranos, cerca de \qty{0.57}{Ga}.
\end{solution}

\begin{solution}{exo:b2:precambrian-earth:5}
(a) $\tfrac14\times1361\times0.706 = \qty{240.2}{W.m^{-2}}$ absorvidos e $\sigma\times288^4 =
\qty{390.1}{W.m^{-2}}$, de modo que $\varepsilon = 240.2/390.1 = 0.616$. (b) $\tfrac14\times1361\times0.4
= \qty{136.1}{W.m^{-2}}$, $T^4 = 136.1/(0.616\times\num{5.670e-8})$, $T = \qty{249.8}{K}$.
(c) A \qty{275.5}{K} o albedo é $0.6 - 0.306\times12.5/20 = 0.4088$: $R = 340.25\times0.5912 =
\qty{201.2}{W.m^{-2}}$ e $E = 0.616\,\sigma\times275.5^4 = \qty{201.2}{W.m^{-2}}$. (d) $R' =
340.25\times0.306/20 = \qty{5.21}{W.m^{-2}.K^{-1}}$, $E' = 4\varepsilon\sigma T^3 =
\qty{2.92}{W.m^{-2}.K^{-1}}$: $R' > E'$, de modo que o equilíbrio é instável
(\cref{prop:b2:precambrian-earth:ebm}).
\end{solution}

\begin{solution}{exo:b2:precambrian-earth:6}
(a) $1/(1 + 0.4\times2.4/4.567) = 0.826$; $1/(1 + 0.4\times0.65/4.567) = 0.946$. (b) O estado
quente precisa de pelo menos 0.932 da luz solar atual: impossível há \qty{2.4}{Ga}, possível
há \qty{0.65}{Ga}. (c) $0.826/0.946 = 0.873$: \qty{13}{\percent} mais fraco, próximo dos
\qty{10}{\percent} citados, que vêm de outro modelo solar.
\end{solution}

\begin{solution}{exo:b2:precambrian-earth:7}
(a) Sturtiana $716.9 - 658.9 = \qty{58.0}{Myr}$; Marinoana $644.5 - 635.4 = \qty{9.1}{Myr}$;
interglacial $658.9 - 644.5 = \qty{14.4}{Myr}$. (b) $(58.0 + 9.1)/85 = 0.79$: cerca de quatro quintos
do período. (c) Seu início só está delimitado entre 649.9 e \qty{639.0}{Ma}: com esses
limites, a Marinoana durou entre cerca de 3 e \qty{15}{Myr}.
\end{solution}

\begin{solution}{exo:b2:precambrian-earth:8}
(a) $\tau = \frac{0.05}{0.95}\times\frac{44.01}{28.97}\times\frac{\num{5.1e18}}{\num{0.26e12}} =
0.05263\times1.519\times\num{1.962e7} = \qty{1.57}{Myr}$. (b) $0.05/\num{278e-6} = 180$ vezes.
(c) O dobro do tempo, \qty{3.14}{Myr}.
\end{solution}

\begin{solution}{exo:b2:precambrian-earth:9}
(a) As pedras caíram em lama mole e encurvaram suas lâminas para baixo: são clastos caídos (dropstones), liberados por
gelo flutuante. (b) Gelo chegando ao mar a \ang{10} do equador é o vestígio esperado de uma
bola de neve, e muito difícil de obter se só as regiões polares estiverem congeladas. (c) Um \omterm{def:b2:precambrian-earth:snowball}{carbonato de capa}:
ele registra o fim súbito da glaciação, com intemperismo intenso e a inundação das
plataformas pela água de degelo.
\end{solution}

\begin{solution}{exo:b2:precambrian-earth:10}
(a) $0.94 > 0.932$, o ponto de virada quente: o equilíbrio quente, $288\times0.94^{1/4} =
\qty{283.6}{K}$, fica logo acima de \qty{283}{K}; e $0.94 < 1.227$, de modo que a bola de neve também existe.
(b) $249.8\times0.94^{1/4} = \qty{246.0}{K}$. (c) $\varepsilon = 0.94\times1361\times0.4/(4\sigma
\times263^4) = 511.7/1085.1 = 0.472$. (d) $\sigma\times263^4 = \qty{271.3}{W.m^{-2}}$, e $(0.616
- 0.472)\times271.3 = \qty{39}{W.m^{-2}}$.
\end{solution}

\begin{solution}{exo:b2:precambrian-earth:11}
(a) $S_{\text{m}}/S_{\text{f}} = (T_1/T_2)^4\,(1 - a_{\text{w}})/(1 - a_{\text{i}})$: a
emissividade se cancela. $(263/283)^4\times0.706/0.4 = 0.7459\times1.765 = 1.317$. (b) Um contraste
menor de albedo entre o gelo e a água aberta, como um gelo mais escuro, empoeirado, ou uma rampa de
temperatura mais larga. (c) Uma bola de neve reflete tanta luz solar que continua fria sob qualquer luz solar ou
efeito estufa com que uma Terra quente congelaria; para degelá-la, o forçamento deve exceder o do
congelamento por toda a largura do laço.
\end{solution}

\begin{solution}{exo:b2:precambrian-earth:12}
Perto de $T^*$, $C\dot\delta = -(E' - R')\delta$, de modo que $\delta \propto e^{-t/\tau}$ com $\tau =
C/(E' - R')$. No equilíbrio quente $R' = 0$ e $E' = 4\times0.616\times\num{5.670e-8}\times
288^3 = \qty{3.34}{W.m^{-2}.K^{-1}}$; $C = 1025\times3992\times50 = \qty{2.05e8}{J.m^{-2}.K^{-1}}$,
portanto $\tau = \qty{6.1e7}{s}$, cerca de \qty{1.9}{yr}. Neste modelo o equilíbrio quente mantém $R' =
0$ até atingir \qty{283}{K}, de modo que $\tau$ fica perto de dois anos até o ponto de virada, e então
o estado desaparece. Com um albedo que varia suavemente, $R'$ cresce à medida que o equilíbrio se aproxima da
dobra, onde $E' = R'$: $\tau$ cresce sem limite, uma desaceleração que anuncia o ponto
de virada.
\end{solution}

\begin{solution}{pb:b2:precambrian-earth:1}
\textbf{1.} $\num{0.0203e-6}\times0.9929\times\num{95.13e-6} = \qty{1.92}{pW.kg^{-1}}$
(\ce{^{238}U}); $\num{0.0203e-6}\times0.00711\times\num{568.47e-6} = \qty{0.0820}{pW.kg^{-1}}$
(\ce{^{235}U}); $\num{0.0795e-6}\times\num{26.28e-6} = \qty{2.09}{pW.kg^{-1}}$ (\ce{^{232}Th});
$\num{240e-6}\times\num{1.20e-4}\times\num{28.47e-6} = \qty{0.820}{pW.kg^{-1}}$ (\ce{^{40}K}).
\textbf{2.} \qty{4.91}{pW.kg^{-1}}.
\textbf{3.} $\ln2/T$: 0.155, 0.985, 0.0495 e \qty{0.555}{Gyr^{-1}}.
\textbf{4.} Fatores 1.52, 14.3, 1.14 e 4.48: 2.91, 1.17, 2.39 e \qty{3.67}{pW.kg^{-1}}.
\textbf{5.} \qty{10.1}{pW.kg^{-1}}, 2.07 vezes o atual.
\textbf{6.} O potássio-40. Sua meia-vida, \qty{1.25}{Gyr}, é curta: mais de duas meias-vidas
se passaram desde então, e mais de três quartos do \ce{^{40}K} daquela época decaíram.
\textbf{7.} $1/(1 + 0.4\times2.7/4.567) = 0.809$; $0.809\times1361 = \qty{1101}{W.m^{-2}}$.
\textbf{8.} $288\times0.809^{1/4} = \qty{273}{K}$, abaixo de \qty{283}{K}: não é um estado quente válido;
a Terra congelaria.
\textbf{9.} $(1101\times0.4/(4\times0.616\,\sigma))^{1/4} = \qty{237}{K}$.
\textbf{10.} $0.616\times0.809 = 0.498$.
\textbf{11.} $0.004/\num{278e-6} = 14$ a $0.03/\num{278e-6} = 108$ vezes.
\textbf{12.} O termostato do intemperismo dos silicatos: o intemperismo desacelera num clima frio, deixando
o dióxido de carbono vulcânico se acumular, e acelera num clima quente.
\textbf{13.} 0.826 da atual, abaixo de 0.932: não há estado quente; a Terra com o ar atual seria uma
bola de neve.
\textbf{14.} $\frac{0.1}{0.9}\times1.519\times\num{1.962e7} = \qty{3.31}{Myr}$.
\textbf{15.} $\frac{0.2}{0.8}\times1.519\times\num{1.962e7} = \qty{7.45}{Myr}$.
\textbf{16.} O \omterm{def:b2:precambrian-earth:snowball}{carbonato de capa}, um dolomito laminado que recobre bruscamente os depósitos glaciais,
precipitado de um mar alimentado por intemperismo intenso no rescaldo quente e inundado pela
água de degelo.
\textbf{17.} O ferro(II) só fica dissolvido em água sem oxigênio: o ferro se acumula no
oceano anóxico sob o gelo e precipita como ferro(III) onde o oxigênio chega, por exemplo com a
água de degelo.
\textbf{18.} $(1616 - 1450)/(32.2 - 23) = \qty{18.0}{\celsius}$ por ponto percentual.
\textbf{19.} $1616 + 18.0\times(30 - 32.2) = \qty{1576}{\celsius}$.
\textbf{20.} \qty{1616}{\celsius}, o ponto medido.
\textbf{21.} $1616 - 18.0\times20.2 = \qty{1252}{\celsius}$.
\textbf{22.} $1616 - 20\times20.2 = \qty{1212}{\celsius}$: \qty{40}{\celsius} mais baixo, um décimo
da diferença entre as lavas; a conclusão não depende da escolha.
\textbf{23.} Os \omterm{def:b2:precambrian-earth:greenstone}{komatiitos} vieram de plumas, mais quentes que o manto ambiente, e de altos graus de
fusão em profundidade; a temperatura de uma lava não é a temperatura potencial de sua fonte. O
manto como um todo era mais quente em 200 a \qty{300}{K}.
\textbf{24.} A lava komatiítica mais quente estava cerca de \textbf{\qty{364}{\celsius}} mais quente que um
basalto primário das dorsais atuais.
\end{solution}
